%
% 6.006 problem set 0 solutions template
%
\documentclass[12pt,twoside]{article}

\input{macros-sp20}
\newcommand{\theproblemsetnum}{0}

\title{6.006 Problem Set 0}

\begin{document}

\handout{Problem Set \theproblemsetnum}{}

\setlength{\parindent}{0pt}
\medskip\hrulefill\medskip

{\bf Name:} HYUNJUN JEON (Luciano)

\medskip\hrulefill

%%%%%%%%%%%%%%%%%%%%%%%%%%%%%%%%%%%%%%%%%%%%%%%%%%%%%
% See below for common and useful latex constructs. %
%%%%%%%%%%%%%%%%%%%%%%%%%%%%%%%%%%%%%%%%%%%%%%%%%%%%%

% Some useful commands:
% $f(x) = \Theta(x)$
% $T(x, y) \leq \log(x) + 2^y + \binom{2n}{n}$
% \ttt{code\_function}


% You can create unnumbered lists as follows:
% \begin{itemize}
%     \item First item in a list
%         \begin{itemize}
%             \item First item in a list
%                 \begin{itemize}
%                     \item First item in a list
%                     \item Second item in a list
%                 \end{itemize}
%             \item Second item in a list
%         \end{itemize}
%     \item Second item in a list
% \end{itemize}

% You can create numbered lists as follows:
% \begin{enumerate}
%     \item First item in a list
%     \item Second item in a list
%     \item Third item in a list
% \end{enumerate}

% You can write aligned equations as follows:
% \begin{align}
%     \begin{split}
%         (x+y)^3 &= (x+y)^2(x+y) \\
%                 &= (x^2+2xy+y^2)(x+y) \\
%                 &= (x^3+2x^2y+xy^2) + (x^2y+2xy^2+y^3) \\
%                 &= x^3+3x^2y+3xy^2+y^3
%     \end{split}
% \end{align}

% You can create grids/matrices as follows:
% \begin{align}
%     A =
%     \begin{bmatrix}
%         A_{11} & A_{21} \\
%         A_{21} & A_{22}
%     \end{bmatrix}
% \end{align}

\begin{problems}

\problem  % Problem 1

\begin{problemparts}
\problempart % Problem 1a
$10$, $15$
\problempart % Problem 1b
$7$
\problempart % Problem 1c
$3$
\end{problemparts}

\problem  % Problem 2

\begin{problemparts}
\problempart % Problem 2a
$\frac{3}{2}$
\problempart % Problem 2b
$\frac{49}{4}$
\problempart % Problem 2c
$\frac{55}{4}$
\end{problemparts}

\problem  % Problem 3

\begin{problemparts}
\problempart % Problem 3a
True
\problempart % Problem 3b
False
\problempart % Problem 3c
True
\end{problemparts}

\problem  % Problem 4

\[
    \sum_{i=1}^{n} i^3 = \left[\frac{n(n+1)}{2}\right]^2.
\]

\textbf{Step 1: Base case ($n=1$).}
\[
    \sum_{i=1}^{1} i^3 = 1
    = \left[\frac{1(1+1)}{2}\right]^2.
\]
Thus, the statement holds for $n=1$.

\textbf{Step 2: Inductive step.}

Assume that the statement holds for $n=k$; that is,
\[
    \sum_{i=1}^{k} i^3 = \left[\frac{k(k+1)}{2}\right]^2.
\]
Then, for $n=k+1$,
\begin{align*}
    \sum_{i=1}^{k+1} i^3
    &= \sum_{i=1}^{k} i^3 + (k+1)^3 \\
    &= \left[\frac{k(k+1)}{2}\right]^2 + (k+1)^3 \\
    &= (k+1)^2\left(\frac{k^2}{4}+k+1\right) \\
    &= (k+1)^2\left(\frac{k^2+4k+4}{4}\right) \\
    &= (k+1)^2\left(\frac{(k+2)^2}{4}\right) \\
    &= \left[\frac{(k+1)(k+2)}{2}\right]^2.
\end{align*}
 \qed
\newpage
\problem  % Problem 5

Proof by induction on $|V|=n$.

\textbf{Step 1: Base case ($n=1$).}

If $|V|=1$, then $|E|=0$. A graph with no edges has no cycle, so the base
case is true.

\textbf{Step 2: Inductive step.}

Assume that a connected graph with $k$ vertices and $k-1$ edges is
acyclic. Now let $G$ have $k+1$ vertices and $k$ edges.

By the handshaking lemma,
\[
    \sum_{v\in V}\deg(v)=2|E|=2k.
\]
Since $G$ is connected, every vertex has degree at least $1$. If every
vertex had degree at least $2$, the degree sum would be at least
$2(k+1)>2k$. Therefore, some vertex $v$ has degree $1$.

Remove $v$ and its edge. The remaining graph has $k$ vertices and $k-1$
edges, and it is still connected. By the inductive hypothesis, it has no
cycle. Adding $v$ back cannot make a cycle because $v$ has only one edge.
Therefore, $G$ is acyclic. \qed

\problem  % Problem 6
Submit your implementation to {\small\url{alg.mit.edu}}.

\begin{lstlisting}
def count_long_subarray(A):
    '''
    Input:  A     | Python Tuple of positive integers
    Output: count | number of longest increasing subarrays of A
    '''
    count = 0
    ##################
    # YOUR CODE HERE #
    ##################
    return count
\end{lstlisting}

\end{problems}

\end{document}
